% AUTO-GENERATED by paper/dai2027/make_tables.py -- do not hand-edit.
% The two listings are lifted verbatim from paper/tables/appendix_prompts.tex,
% which paper/analysis/make_prompt_appendix.py renders from the live prompts.

This appendix and the schema of Appendix~\ref{sec:appendix-schema} are what this paper discloses in place of a code release (see the Ethical Statement). The Pass-1 prompt and the classification rule are reproduced verbatim, as the system emits them with the document content replaced by a placeholder; the other rules the body cites are summarized. \textbf{The complete text of all three prompts is provided as supplementary material (\texttt{prompts.txt}).}

\subsection{Pass 1: depth-map inference}
\label{sec:prompt-depthmap}
Run once per document, on a layout-stratified sample of text blocks (\S\ref{sec:method-detection}). Its output is supplied to every detection call as the authority for what each nesting level is.

\begin{lstlisting}[style=promptstyle]
You are analyzing the structural skeleton of an early learning standards document. You will NOT extract individual elements — only the document's nesting hierarchy.

Read the sample below and identify how many distinct nesting depths the document uses, what each depth looks like (its prefix/label pattern), and one concrete example. The DEEPEST depth is always the leaf (individual learning goals/foundations/benchmarks).

This document's depths must then be mapped to a canonical 4-level hierarchy:
- depth 1 → "domain"
- depth 2 → "strand"
- depth 3 → "sub_strand"
- deepest depth → "indicator"

Rules:
- Fill the canonical levels FROM THE TOP DOWN. The first depth is always "domain", the deepest depth is always "indicator", and any depth in between takes the next canonical level in order: the depth directly under domain is "strand", the one below that is "sub_strand". A document with fewer depths therefore drops the levels at the BOTTOM of that middle range — never "strand".
- A 4-level document (domain > group > sub-group > leaf) maps to: domain, strand, sub_strand, indicator.
- A 3-level document (domain > group > leaf) maps to: domain, strand, indicator. There is NO sub_strand — the single middle level is a STRAND, never a sub_strand.
- A 2-level document (domain > leaf) maps to: domain, indicator. There is no strand and no sub_strand.
- Use depth POSITION, not the document's labels. If a document uses "Sub-Strand" as the label for the second level (directly under domain), it is still a STRAND in our canonical hierarchy — and likewise a document that calls its single middle level "Section", "Concept", "Goal", or anything else still maps that level to "strand".

Every line is prefixed with its page and the horizontal position it was laid out at: `[Page N | x=0.09]`, where x is the line's left edge as a fraction of the page width. Use x to work out whether a page uses a side-by-side layout: lines of one column share nearly the same x, and a jump to a clearly different x means a different column. Lines you judge to be in the same column read as one continuous cell, even when lines from other columns are interleaved between them. Treat parallel columns as repetitions of ONE depth, not as separate depths.

Be deterministic and conservative. Do not speculate. If you cannot tell whether two depths are distinct, assume they are the same depth.

Output ONLY a JSON object with this exact shape:
{
  "doc_depths": [
    {"depth": 1, "canonical_level": "domain",     "label_in_doc": "<what the doc calls this>", "prefix_pattern": "<regex-ish pattern e.g. 'ALL-CAPS HEADING' or 'N. <Title>: <desc>'>", "example": "<exact text from doc>"},
    {"depth": 2, "canonical_level": "strand|indicator", ...}
  ],
  "notes": "<one short sentence about the LEVEL LAYOUT only, e.g. 'age-band columns repeat one depth'. A later stage reads this as an instruction, so: never describe text below the deepest depth as examples/illustrations/anecdotes, and never say to ignore or skip anything. Empty string if nothing unusual.>"
}

DOCUMENT SAMPLE:

<<serialized document lines, one per row, each prefixed
  [Page N | x=<left edge as fraction of page width>]>>
\end{lstlisting}

\subsection{Pass 2: per-chunk detection}
\label{sec:prompt-detection}
The prompt opens with the canonical hierarchy, the layout-coordinate convention and the classification rule, then states nine numbered extraction rules and the output schema. Rule 1 requires every element to be backed by a heading line that literally appears in the chunk, so a chunk that opens part-way through a section emits the children alone and the parent is recovered in parsing.

\subsubsection{Classification is positional}
\label{rule:positional}
\begin{lstlisting}[style=promptstyle]
CLASSIFICATION RULE:
- If a depth map is provided, look up each element's depth in the document and use the `canonical_level` from that depth. Do NOT reclassify based on prefix style.
- If no depth map is provided, classify by nesting POSITION, never by what the document calls a level.
\end{lstlisting}

\subsubsection{Code resolution (rule 4)}
\label{rule:code-resolution}
The longest rule. \texttt{code} is required and never \texttt{null}. A printed code always wins, and the rule orders three places to look for one: (i)~inline on the heading line, including the \texttt{<Label> <id>} form for any label word; (ii)~a caption beside the heading, such as a parenthetical or column header, copied verbatim; (iii)~the shared leading prefix of the element's descendants' dotted codes, peeled one whole segment per level, agreed by at least two descendants, and never taken from a code that carries a structural label word. Only when all three fail is a code of at most five characters derived from the title: split on spaces and slashes, drop a fixed list of connector words, then take the first four letters of a single remaining word or the initials of up to five. That derivation is also executed in Python (\S\ref{sec:lessons}, lesson~1).

\subsection{Parsing}
\label{sec:prompt-parsing}
The parsing prompt states the output schema, then rules for composing each indicator's fully qualified code from its chain of ancestors. A code that is already a dotted path beginning with its own domain's code is used as is and never re-prefixed, at every level. A \texttt{<Label> <id>} code (\texttt{Benchmark 1.1}) contributes only its identifier, for any label word. Peeling ancestors off a descendant's code stops where the printed namespace stops, so a level outside it builds its code from its own heading's identifier (Nevada's strand \texttt{SS.1} above sub-strand \texttt{SS.ID}). Side-by-side columns are disambiguated by a suffix on the indicator code only: a month range for age columns (\texttt{.36-48}), otherwise the column label's first four letters (\texttt{.DISC}). The shape repairs of \S\ref{sec:method-parsing} that collapse a duplicated parent segment or decline an unverifiable repair are deterministic Python applied to the model's output, not prompt text.
